\documentclass[conference]{IEEEtran}
\IEEEoverridecommandlockouts
\usepackage{cite}
\usepackage{amsmath,amssymb,amsfonts}
\usepackage{algorithmic}
\usepackage{graphicx}
\usepackage{textcomp}
\usepackage{xcolor}
\usepackage{booktabs}
\usepackage{multirow}
\usepackage{hyperref}
\def\BibTeX{{\rm B\kern-.05em{\sc i\kern-.025em b}\kern-.08em
    T\kern-.1667em\lower.7ex\hbox{E}\kern-.125emX}}
\begin{document}

\title{Audio-Augmented Diffusion Policy: Sound as Pseudo-Instrumentation for Contact-Rich Robot Manipulation}

\author{\IEEEauthorblockN{1\textsuperscript{st} Given Name Surname}
\IEEEauthorblockA{\textit{dept. name of organization (of Aff.)} \\
\textit{name of organization (of Aff.)}\\
City, Country \\
email address or ORCID}
\and
\IEEEauthorblockN{2\textsuperscript{nd} Given Name Surname}
\IEEEauthorblockA{\textit{dept. name of organization (of Aff.)} \\
\textit{name of organization (of Aff.)}\\
City, Country \\
email address or ORCID}
}

\maketitle

\begin{abstract}
Contact-rich manipulation benefits from object-mounted instrumentation: a contact switch on a button, for example, enables automated expert data collection via a scripted proportional controller and provides a direct task-progress signal during policy execution.
But instrumenting every object is impractical outside the lab.
We investigate whether audio can partially substitute for instrumentation at test time while retaining its data-collection advantages.
We integrate an Audio Spectrogram Transformer (AST) into the conditioning pathway of a diffusion policy and compare frozen versus finetuned backbones, embedding versus classifier outputs, on a real UR3e button-press-and-retract task with 40 automatically collected demonstrations.
Our best audio-augmented model achieves 57.5\% success, compared to 55.0\% for a vision-only baseline and 62.5\% for an oracle with direct button-state access.
Task-specific AST finetuning is critical: generic AudioSet features degrade performance below the vision-only baseline, while a button-click-finetuned AST recovers part of the gap to the oracle.
\end{abstract}

\begin{IEEEkeywords}
imitation learning, diffusion policy, audio, multimodal, contact-rich manipulation
\end{IEEEkeywords}

%=============================================================================
\section{Introduction}
%=============================================================================

Learning manipulation policies from demonstrations has seen rapid progress with diffusion-based approaches~\cite{chi2023diffusion}.
Contact-rich tasks, however, remain difficult for vision-only systems.
In a button-pressing task, the robot must detect the moment of successful contact to stop pressing and retract.
This event is often visually ambiguous: the end-effector occludes the contact point, and the physical deflection is sub-millimeter.

Object-mounted instrumentation, such as a binary contact switch, solves this directly.
In our setting, a contact switch on the button provides three distinct advantages:
\begin{enumerate}
    \item \emph{Automated data collection.}
    The contact signal enables a scripted proportional controller that collects demonstrations without human teleoperation.
    The controller uses the button signal to detect successful presses and trigger retraction automatically.
    \item \emph{Demonstration quality.}
    The scripted controller produces consistent, minimal-force demonstrations.
    It approaches with proportional speed control and stops immediately upon contact, avoiding the excess force and timing variability typical of human teleoperation.
    \item \emph{Privileged task-progress signal.}
    Conditioning a learned policy on the button state during execution improves performance by providing direct knowledge of whether contact has occurred.
\end{enumerate}

The first two advantages apply only during data collection and are preserved regardless of what the policy observes at deployment.
The third requires the instrument at test time.
Our question is whether this privileged signal can be replaced by a cheaper modality: the sound of a button click, captured by an ambient microphone, requires no modification to the object or environment.

We investigate this by integrating a pretrained Audio Spectrogram Transformer (AST)~\cite{gong2021ast} into the global conditioning pathway of a diffusion policy~\cite{chi2023diffusion}.
The AST processes a mel-spectrogram from the last 3~seconds of ambient audio, and its features are concatenated with visual features from a ResNet-18 encoder before conditioning the denoising U-Net.
We compare:
\begin{itemize}
    \item Vision only: RGB wrist camera + proprioceptive state.
    \item Oracle: vision + ground-truth binary button signal.
    \item Audio embeddings: vision + frozen or finetuned AST pooler output.
    \item Audio classifier (pseudo-instrumentation): vision + AST logits from a button-click classifier, replacing the ground-truth button signal.
\end{itemize}

On a real UR3e robot, the best audio-only model reaches 57.5\% success compared to 55.0\% for vision only and 62.5\% for the oracle.
Generic AudioSet features consistently hurt performance; task-specific finetuning is required for audio to help.

%=============================================================================
\section{Related Work}
%=============================================================================

\textbf{Diffusion policies for manipulation.}
Chi et al.~\cite{chi2023diffusion} introduced Diffusion Policy, framing visuomotor policy learning as conditional denoising diffusion over action trajectories.
A 1D U-Net conditioned on visual features (ResNet with spatial softmax) and proprioceptive state uses DDPM~\cite{ho2020ddpm} to generate multi-step action chunks.
Follow-up work has extended diffusion policies to 3D representations~\cite{ze20243ddp}, language conditioning~\cite{ha2023scaling}, and dexterous manipulation~\cite{chi2024universal}.
We extend the conditioning pathway to include audio features.

\textbf{Audio in robotics.}
Audio has been used for object recognition~\cite{gandhi2020swoosh}, material classification~\cite{clarke2018learning}, and grasp stability estimation~\cite{huang2024audiograsp}.
Li et al.~\cite{li2022seehearfeel} showed that audio and tactile signals improve policy performance on contact-rich tasks through a self-supervised multimodal encoder.
Liu et al.~\cite{liu2024maniwav} integrate audio into diffusion policies using a contrastive audio-visual representation.
Our approach differs in directly conditioning the diffusion policy on a pretrained AST, comparing both generic embeddings and task-specific classifier outputs.

\textbf{Audio Spectrogram Transformer.}
Gong et al.~\cite{gong2021ast} proposed the AST, a pure-attention model that treats audio spectrograms as image-like inputs to a Vision Transformer.
Pretrained on AudioSet~\cite{gemmeke2017audioset}, the AST achieves strong transfer performance on downstream audio tasks.
We use this pretrained representation and investigate whether AudioSet features transfer to robotic contact detection.

%=============================================================================
\section{Method}
%=============================================================================

\subsection{Task and Hardware}

The task is \emph{button press and retract}: the robot lowers its end-effector onto a push-button until the button clicks, then retracts upward.
Success requires reaching the button, pressing with sufficient but not excessive force, and detecting the click to initiate retraction.

We use a UR3e robot arm with a Robotiq 2F-85 gripper, a wrist-mounted Intel RealSense D435 RGB camera ($320{\times}240$ at 30~Hz), a Bota Systems force-torque sensor, and an ambient microphone.
The button is instrumented with a binary contact switch used for ground-truth labeling during data collection and as oracle input during policy evaluation.
Control runs at 10~Hz.

\subsection{Automated Data Collection}
\label{sec:data_collection}

The contact switch enables fully automated demonstration collection through a three-phase proportional controller, eliminating the need for human teleoperation.

\textbf{Phase~1 (lateral alignment).}
The controller computes the 3D vector from end-effector to button in the tool frame and applies proportional lateral corrections with gain $k_{xy}$.
This phase terminates when lateral error drops below 2~mm.

\textbf{Phase~2 (descent).}
The controller descends toward the button with proportional vertical gain $k_z$.
Below a pre-contact height, descent slows to a constant 1~mm/step.
This phase terminates when the contact switch detects a button press.

\textbf{Phase~3 (retraction).}
The controller moves upward at a fixed 2~cm/step until reaching 10~cm above the button.

To produce diverse demonstrations, we randomize both the initial pose ($\pm$5~cm in XY, 0--5~cm in Z, $\pm$90$^\circ$ rotation around Z) and the approach style.
An approach-style parameter $\alpha$ is sampled log-uniformly from $[0.1, 10]$ and modulates the gain ratio: $k_{xy} = k_0\sqrt{\alpha}$, $k_z = k_0/\sqrt{\alpha}$.
Low $\alpha$ produces aggressive vertical approaches; high $\alpha$ produces careful lateral alignment before descent.
Step sizes are clipped to $[2\text{~mm}, 3\text{~cm}]$ for safety.
Controller hyperparameters were tuned in a 2D simulation prior to deployment.
We collected 40 demonstrations with this system.

\subsection{Observation and Action Spaces}

\textbf{Observations.}
Each observation at time $t$ consists of:
\begin{itemize}
    \item Wrist RGB image $\mathbf{I}_t \in \mathbb{R}^{3 \times 224 \times 224}$, resized from the raw camera feed.
    \item Robot state $\mathbf{s}_t$: a scalar button signal (in oracle and combined models) or empty (in vision-only and audio-only variants).
    \item Audio spectrogram $\mathbf{A}_t \in \mathbb{R}^{298 \times 128}$: a mel-spectrogram computed over the most recent 3~seconds of audio via Kaldi, yielding 298 time frames $\times$ 128 frequency bins.
\end{itemize}

\textbf{Actions.}
Actions are 3D Cartesian position deltas $\mathbf{a}_t = [\Delta x, \Delta y, \Delta z]$ expressed in the tool frame: $\mathbf{a}_t = \mathbf{R}_t^{-1}(\mathbf{p}_{\text{target}} - \mathbf{p}_t)$, where $\mathbf{R}_t$ and $\mathbf{p}_t$ are the end-effector orientation and position.
This frame choice makes actions invariant to the robot's base-frame orientation.

\subsection{Diffusion Policy Architecture}

We build on Diffusion Policy~\cite{chi2023diffusion}, which models the action distribution as a denoising diffusion process.
A 1D convolutional U-Net with FiLM conditioning~\cite{perez2018film} generates a horizon of $H{=}16$ future actions, of which 8 are executed before re-planning.
The U-Net is conditioned on a global feature vector assembled from the $n{=}2$ most recent observations:

\begin{equation}
    \mathbf{g} = \text{concat}\!\left[\,\mathbf{s}_{t-1:t},\; \mathbf{f}^{\text{rgb}}_{t-1:t},\; \mathbf{f}^{\text{audio}}_{t-1:t}\,\right]
    \label{eq:global_cond}
\end{equation}

\noindent where $\mathbf{s}$ is the robot state, $\mathbf{f}^{\text{rgb}}$ are visual features, and $\mathbf{f}^{\text{audio}}$ are audio features (omitted in the vision-only baseline).

\subsubsection{Vision Encoder}
A ResNet-18 backbone with group normalization extracts convolutional feature maps from each wrist image.
A spatial softmax layer~\cite{finn2016spatial} computes 32 expected keypoint positions, yielding a 64-dimensional vector per image projected to 32 dimensions.
Random crops ($84{\times}84$) are applied during training; center crops during evaluation.

\subsubsection{Audio Encoder}
\label{sec:audio_encoder}
We use an AST~\cite{gong2021ast} pretrained on AudioSet~\cite{gemmeke2017audioset}.
The model treats the $298 \times 128$ mel-spectrogram as a sequence of patches and produces a 768-dimensional CLS token embedding.
Spectrograms are normalized using training-set statistics: $\hat{\mathbf{A}} = (\mathbf{A} - \mu) / (2\sigma)$.

We explore several integration strategies:
\begin{itemize}
    \item Frozen embedding: the AST backbone is frozen; the 768-dim pooler output is concatenated into the global conditioning vector.
    \item Unfrozen embedding: the full AST is finetuned end-to-end with the policy.
    \item Frozen/unfrozen classifier: the AST is first finetuned as a binary button-click detector (2-class classification, lr $10^{-5}$, 10 epochs). Its 2-dimensional logits then replace the ground-truth button signal in the policy's state input. This is our \emph{pseudo-instrumentation} approach.
\end{itemize}

\subsubsection{Denoising U-Net}
The global conditioning vector is injected into each residual block of a 1D convolutional U-Net via FiLM layers.
The U-Net has channel dimensions (512, 1024, 2048), kernel size~5, and group normalization with 8 groups.
The diffusion process uses DDPM with 100 forward steps, a squared-cosine beta schedule, and epsilon prediction.

\subsection{Training}

All models are trained with Adam (lr $10^{-4}$, weight decay $10^{-6}$), a cosine schedule with 500 warmup steps.
Image augmentations (brightness, contrast, saturation, hue, sharpness jitter) are applied.
Actions and states are normalized to $[-1, 1]$ via min-max scaling; spectrograms use AST-specific normalization (Section~\ref{sec:audio_encoder}).

%=============================================================================
\section{Experiments}
%=============================================================================

\subsection{Evaluation Protocol}

Each model is evaluated over $N{=}40$ real-robot episodes.
We classify each episode into one of five outcome categories:
\begin{itemize}
    \item \textbf{Success}: the button is pressed and the robot retracts.
    \item \textbf{Hover}: the robot positions above the button but fails to make sufficient contact.
    \item \textbf{Too hard}: the robot presses with excessive force, triggering safety limits.
    \item \textbf{Sideways}: the gripper approaches from an incorrect angle.
    \item \textbf{Fake}: the policy retracts before the button is actually pressed (audio false positive).
\end{itemize}

Force-torque data from the wrist sensor is recorded for post-hoc contact analysis.

\subsection{Results}

Table~\ref{tab:main_results} presents results for all 12 model variants.

\begin{table}[t]
\caption{Results on button-press-and-retract ($N{=}40$ per model). Succ = successes, Hov = hover failures, TH = too hard, SW = sideways, FK = fake retraction.}
\label{tab:main_results}
\begin{center}
\setlength{\tabcolsep}{3pt}
\begin{tabular}{lcccccc}
\toprule
\textbf{Model} & \textbf{Rate} & \textbf{Succ} & \textbf{Hov} & \textbf{TH} & \textbf{SW} & \textbf{FK} \\
\midrule
Expert (scripted) & 100\% & 40 & 0 & 0 & 0 & 0 \\
\midrule
Vision only & 55.0\% & 22 & 16 & 0 & 2 & 0 \\
Oracle (V + button) & 62.5\% & 25 & 15 & 0 & 0 & 0 \\
\midrule
\multicolumn{7}{l}{\textit{Audio + button signal}} \\
\quad + GenAST frozen & 45.0\% & 18 & 21 & 0 & 1 & 0 \\
\quad + GenAST unfrozen & 62.5\% & 25 & 15 & 0 & 0 & 0 \\
\midrule
\multicolumn{7}{l}{\textit{Generic AST (no button finetuning)}} \\
\quad Frozen embed & 47.5\% & 19 & 21 & 0 & 0 & 0 \\
\quad Unfrozen embed & 40.0\% & 16 & 24 & 0 & 0 & 0 \\
\midrule
\multicolumn{7}{l}{\textit{Button-finetuned AST}} \\
\quad Pseudo-instr (swap) & 55.0\% & 22 & 17 & 0 & 0 & 1 \\
\quad Frozen classif & \textbf{57.5\%} & 23 & 17 & 0 & 0 & 0 \\
\quad Unfrozen classif & 42.5\% & 17 & 23 & 0 & 0 & 0 \\
\quad Frozen embed & 45.0\% & 18 & 14 & 3 & 1 & 4 \\
\quad Unfrozen embed & \textbf{57.5\%} & 23 & 17 & 0 & 0 & 0 \\
\bottomrule
\end{tabular}
\end{center}
\end{table}

\subsubsection{Baselines}

The vision-only model achieves 55.0\% success.
Its failures are predominantly hovering (16/40): the robot positions itself above the button but does not descend far enough to make contact.
Two episodes fail due to approaching from the side.
No episodes involve excessive force.

The oracle, which receives the ground-truth button signal, reaches 62.5\%.
All 15 failures are hovers.
The button signal helps the robot know when to retract but does not help it reach the button in the first place.

\subsubsection{Audio-only models}
The best audio-only configurations are the button-finetuned AST models.
The frozen classifier achieves 57.5\%, as does the unfrozen embedding, gaining 2.5 percentage points over the vision-only baseline and closing about a third of the 7.5-point gap to the oracle.
The pseudo-instrumentation variant (AST logits replacing the button signal) matches the vision-only baseline at 55.0\%.

Generic AST features, without button-click finetuning, hurt performance.
The frozen generic model reaches only 47.5\%, and the unfrozen model drops to 40.0\% with 24/40 hover failures.
AudioSet features, trained on 527 general sound categories, do not transfer to this contact-detection task without additional finetuning.

\subsubsection{Combined audio and button}
Adding an unfrozen generic AST to the oracle model matches the oracle at 62.5\%.
The frozen variant drops to 45.0\%, suggesting that untuned audio features can interfere with the button signal.

\subsubsection{Failure modes}
Hovering dominates across all models (14--24 of 40 episodes in every variant except the expert), indicating that the primary bottleneck is spatial: reaching the button from a randomized initial pose.
Neither audio nor the button signal reduces hovering, because both provide information about contact that only becomes relevant once the robot is near the button.

The ``fake'' failure mode occurs when the AST signals a button press that has not happened, causing premature retraction.
This affects the frozen embedding model most (4/40) and appears once in the pseudo-instrumentation model.
It is absent from all non-audio models.

The ``too hard'' mode appears only in the button-finetuned frozen embedding model (3/40), which also has the most fake retractions.
This model's diverse failure profile suggests unstable conditioning, possibly because the 768-dimensional frozen embedding introduces too much noise relative to its signal.

\subsubsection{Intermediate projection layer}

Table~\ref{tab:intermediate} tests whether a 64-dimensional bottleneck between the AST and the conditioning vector helps by filtering noise from the high-dimensional embedding.

\begin{table}[t]
\caption{Effect of a 64-dim intermediate projection on generic AST models ($N{=}40$).}
\label{tab:intermediate}
\begin{center}
\begin{tabular}{lccc}
\toprule
\textbf{Model} & \textbf{Rate} & \textbf{Hov} & \textbf{SW} \\
\midrule
Vision only & 55.0\% & 16 & 2 \\
\midrule
GenAST frozen & 47.5\% & 21 & 0 \\
GenAST frozen + 64-dim & 42.5\% & 22 & 1 \\
GenAST unfrozen & 40.0\% & 24 & 0 \\
GenAST unfrozen + 64-dim & 40.0\% & 24 & 0 \\
\bottomrule
\end{tabular}
\end{center}
\end{table}

It does not.
For the frozen model, adding the projection slightly hurts performance (47.5\% to 42.5\%); for the unfrozen models, results are unchanged.
When the underlying features are not task-relevant, compressing them through a bottleneck does not help.

\subsubsection{Force analysis}

We analyze peak vertical force ($F_z$) during contact episodes, excluding ``too hard'' and ``sideways'' episodes to focus on successful and hovering contact profiles.
We apply Mann-Whitney U pairwise tests with Bonferroni correction ($p > 0.0008$ threshold for 78 comparisons) and report compact letter display (CLD) groupings.

The scripted expert produces the lightest contact forces (CLD group~\textsf{a}).
The oracle falls in group~\textsf{ab}, pressing more gently than the vision-only baseline (group~\textsf{c}).
Several button-finetuned AST models fall in group~\textsf{bc}, intermediate between oracle and vision-only.
Models without contact feedback (vision-only, generic AST, pseudo-instrumentation) cluster in group~\textsf{c}, pressing harder on average.
This pattern is consistent with the role of contact detection: models that can identify the moment of button press exert less force because they initiate retraction sooner.

\subsubsection{Statistical analysis}

With $N{=}40$ episodes per model, the 2.5 percentage point improvement of the best audio models over the vision-only baseline is modest.
Bayesian 95\% credible intervals (Beta$(k{+}1, n{-}k{+}1)$ posterior with uniform prior) overlap between most model variants.
The vision-only baseline's interval spans roughly 40--70\%, and the best audio models' intervals are similarly wide.
The difference between vision-only and oracle (7.5 percentage points) shows clearer separation but still overlaps.
Larger sample sizes would be needed to establish statistical significance for the observed differences.

%=============================================================================
\section{Discussion}
%=============================================================================

\textbf{Instrumentation as a data-collection tool.}
The contact switch's most unambiguous benefit in our pipeline is enabling automated demonstration collection.
The scripted proportional controller avoids the variability and excess force of human teleoperation, producing demonstrations where the expert achieves 100\% success with the lightest contact forces of any model.
These data-collection advantages persist at deployment regardless of whether the policy observes the button signal.

\textbf{Audio provides a small contact signal.}
The best audio-only models (button-finetuned AST, frozen classifier and unfrozen embedding) reach 57.5\%, a 2.5 percentage point gain over vision alone.
This improvement is consistent with audio providing some contact feedback, but the effect is modest relative to the 7.5-point gap between vision-only and oracle.
The force analysis corroborates this: button-finetuned AST models press with intermediate force, falling between the oracle (which presses lightly) and the vision-only baseline (which presses harder on average).

\textbf{Task-specific finetuning is essential.}
Generic AudioSet features do not transfer to robotic contact detection.
Both frozen (47.5\%) and unfrozen (40.0\%) generic AST models perform worse than the vision-only baseline (55.0\%).
Finetuning the AST on button-click spectrograms (a binary classification task requiring only labeled audio segments) is necessary for audio features to help.

\textbf{Hovering is the bottleneck.}
The dominant failure across all models is hovering: the robot stops above the button without making contact.
This is a spatial reaching problem, not a contact-detection problem, and neither audio nor the button signal addresses it.
Improving reaching accuracy through better visual representations, depth sensing, or policy architecture changes would likely yield larger gains than further audio integration.

\textbf{Limitations.}
This study covers a single task with a distinctive acoustic signature.
Button clicks are loud and brief; generalization to quieter or more ambiguous contact sounds (e.g., soft insertion) is unknown.
The 40-episode evaluation, while standard for real-robot experiments, limits statistical power.
The ambient microphone also captures robot motor noise, which may confound results in noisier settings.

%=============================================================================
\section{Conclusion}
%=============================================================================

We integrated an Audio Spectrogram Transformer into a diffusion policy for a real-robot button-press task, testing whether audio can substitute for object-mounted instrumentation at deployment.
Instrumentation proved most valuable for automating data collection: a scripted proportional controller produced 40 expert demonstrations without human teleoperation.
At test time, button-finetuned AST features provided a modest improvement over vision alone (57.5\% vs.\ 55.0\%), closing about a third of the gap to an oracle with direct button-state access (62.5\%).
Generic AudioSet features consistently degraded performance, confirming that task-specific finetuning is necessary.
The primary failure mode across all models was hovering rather than contact detection, suggesting that spatial reaching accuracy is the more pressing bottleneck for this task.

\section*{Acknowledgment}

% TODO: Add acknowledgments.

\begin{thebibliography}{00}
\bibitem{chi2023diffusion} C. Chi, S. Feng, Y. Du, Z. Xu, E. Cousineau, B. Burchfiel, and S. Song, ``Diffusion policy: Visuomotor policy learning via action diffusion,'' in \textit{Proc. Robotics: Science and Systems (RSS)}, 2023.
\bibitem{ho2020ddpm} J. Ho, A. Jain, and P. Abbeel, ``Denoising diffusion probabilistic models,'' in \textit{Advances in Neural Information Processing Systems (NeurIPS)}, 2020.
\bibitem{gong2021ast} Y. Gong, Y.-A. Chung, and J. Glass, ``AST: Audio Spectrogram Transformer,'' in \textit{Proc. Interspeech}, 2021.
\bibitem{gemmeke2017audioset} J. F. Gemmeke, D. P. W. Ellis, D. Freedman, A. Jansen, W. Lawrence, R. C. Moore, M. Plakal, and M. Ritter, ``Audio Set: An ontology and human-labeled dataset for audio events,'' in \textit{Proc. IEEE ICASSP}, 2017.
\bibitem{perez2018film} E. Perez, F. Strub, H. de Vries, V. Dumoulin, and A. Courville, ``FiLM: Visual reasoning with a general conditioning layer,'' in \textit{Proc. AAAI}, 2018.
\bibitem{finn2016spatial} C. Finn, X. Y. Tan, Y. Duan, T. Darrell, S. Levine, and P. Abbeel, ``Deep spatial autoencoders for visuomotor learning,'' in \textit{Proc. IEEE ICRA}, 2016.
\bibitem{ze20243ddp} Y. Ze, G. Zhang, K. Zhang, C. Hu, M. Wang, and H. Xu, ``3D diffusion policy: Generalizable visuomotor policy learning via simple 3D representations,'' in \textit{Proc. RSS}, 2024.
\bibitem{ha2023scaling} H. Ha, P. Florence, and S. Song, ``Scaling up and distilling down: Language-guided robot skill acquisition,'' in \textit{Proc. CoRL}, 2023.
\bibitem{chi2024universal} C. Chi, Z. Xu, C. Pan, E. Cousineau, B. Burchfiel, S. Feng, R. Tedrake, and S. Song, ``Universal manipulation interface: In-the-wild robot teaching without in-the-wild robots,'' in \textit{Proc. RSS}, 2024.
\bibitem{gandhi2020swoosh} D. Gandhi, A. Gupta, and L. Pinto, ``Swoosh! Rattle! Thump! Actions that sound,'' in \textit{Proc. RSS}, 2020.
\bibitem{clarke2018learning} S. Clarke, T. Rhodes, C. G. Atkeson, and O. Kroemer, ``Learning audio feedback for estimating amount and flow of granular material,'' in \textit{Proc. CoRL}, 2018.
\bibitem{huang2024audiograsp} T. Huang, I. Guzey, and S. Chintala, ``Hearing touch: Audio-visual pretraining for contact-rich manipulation,'' in \textit{Proc. ICRA}, 2024.
\bibitem{li2022seehearfeel} J. Li, Y. Dong, R. Calandra, and P. Abbeel, ``See, hear, and feel: Smart sensory fusion for robotic manipulation,'' in \textit{Proc. CoRL}, 2022.
\bibitem{liu2024maniwav} Z. Liu, C. Chi, E. Cousineau, N. Kuppuswamy, B. Burchfiel, and S. Song, ``ManiWAV: Learning robot manipulation from in-the-wild audio-visual data,'' in \textit{Proc. CoRL}, 2024.
\end{thebibliography}

\end{document}
